% Generated draft: requires mathematical and source review.
\documentclass[11pt]{article}
\usepackage[margin=25mm]{geometry}
\usepackage{fontspec}
\setmainfont{Latin Modern Roman}
\usepackage{amsmath,amssymb,mathtools,unicode-math}
\setmathfont{Latin Modern Math}
\usepackage{enumitem,needspace,xurl,hyperref}
\setlength{\parindent}{0pt}
\setlength{\parskip}{5pt}
\setlength{\emergencystretch}{3em}
\newcommand{\N}{\mathbb{N}}
\newcommand{\Z}{\mathbb{Z}}
\newcommand{\Q}{\mathbb{Q}}
\newcommand{\R}{\mathbb{R}}
\newcommand{\C}{\mathbb{C}}
\newcommand{\F}{\mathbb{F}}
\newcommand{\A}{\mathbb{A}}
\newcommand{\PP}{\mathbb{P}}
\newcommand{\E}{\mathbb{E}}
\newcommand{\eps}{\varepsilon}
\newcommand{\dd}{\,\mathrm{d}}
\begin{document}
\section*{UnsolvedMath ID: 6700010}
\textbf{Source problem code:} catalogue entry
\subsection{Definitions and mathematical statement}
The main problem is to extend positivity of scalar curvature from smooth Riemannian manifolds to singular Fano varieties. A smooth complex Fano variety has ample anticanonical bundle \(K_X^{-1}=\det T_X\). In the usual meaning of ampleness, a sufficiently high power gives a projective embedding; the source's footnote merely says a power is generated by sections, which by itself is weaker and needs clarification.

Construct a notion denoted \(\operatorname{Sc}(X)\), along the lines of ``Question 1 of the previous section'', such that the condition
\[
 \operatorname{Sc}(X)\ge0
\]
is preserved under the intended deformations of smooth Fano varieties to singular ones, including low-degree hypersurface degenerations in \(\mathbb{CP}^{m+1}\). The referenced definition proposal, metric and measure data, family/convergence notion, and admissible singularities are absent. In particular a scalar-curvature notation for an algebraic variety needs to specify its metric dependence. No restriction to normal or \(\mathbb Q\)-Fano varieties is silently made: the excerpt explicitly discusses nonnormal, locally reducible behavior.

The source warns that small-ball volumes at singular points can be significantly larger than at regular points, especially in the locally reducible setting, so the naive smooth small-ball-volume definition is not directly applicable. The desired degeneration-stable nonnegative condition is distinguished from merely having positive scalar curvature on the regular locus.

The supplied passage then moves into surgery background. It reports preservation of scalar lower bounds under the relevant surgery construction, connected-sum preservation for positive-scalar-curvature \(n\)-manifolds when \(n\ge3\), and making orientable examples simply connected by attaching 2-handles in dimension at least \(4\). These are attributed contextual claims, not new hypotheses of the Fano problem. The excerpt's wording ``codimension 2 surgery'' sits beside surgery on submanifolds of dimension at most \(n-3\), so the indexing convention needs verification. Its tube discussion uses \(\partial U_\varepsilon(P)\) with sphere fibers of dimension \(k\ge2\) when \(\operatorname{codim}(P)\ge3\), and a blow-up scale \(\varepsilon^{-2}\). A round radius-\(\varepsilon\) \(k\)-sphere has scalar curvature \(k(k-1)\varepsilon^{-2}\); the extracted \(\varepsilon^2\) factor and the notation \(X_\varepsilon=U_\varepsilon(P)\) versus its boundary are ambiguous. The final polyhedral-subset sentence is truncated and supplies no additional complete question.

The summary reports related singular K\"ahler--Einstein, stability, orbifold, and synthetic-Ricci progress, mentioning Donaldson--Sun and singular Yau--Tian--Donaldson theory, while saying no unified definition for the full requested class was verified. Those frameworks are not identified with the proposed generalized scalar curvature without exact hypotheses. The partial-progress label remains an unverified source assessment.
\subsection{Short English statement}
Define positive or nonnegative generalized scalar curvature for singular Fano varieties so that nonnegative curvature survives the intended smooth-to-singular degenerations. The missing proposed definition, metric data, singularity class, and deformation notion must be supplied; the attached surgery discussion also contains extraction and convention issues.
\subsection{Sources}
\begin{itemize}
\item Local registry: \url{data/unsolvedmath/problems.json}, permanent UnsolvedMath ID \(6700010\), source code catalogue entry; complete supplied statement, background, and research summary, updated 2026-08-06.
\item Misha Gromov, \emph{101 Questions, Problems and Conjectures around Scalar Curvature} (2017), printed tag [?10], occurrence 10, PDF page 9: \url{https://www.ihes.fr/~gromov/wp-content/uploads/2018/08/101-problemsOct1-2017.pdf}. Supplied normalized PDF-text extraction accessed 2026-07-29; described as an incomplete source.
\item Source reference [11], Basilio and collaborators, ``Sewing Riemannian Manifolds'' (2017), as extracted; full bibliographic details are not supplied.
\item Source reference [122], Walsh, ``Metrics of positive scalar'' (2008), as extracted; the title is incomplete and full details are not supplied.
\item The summary mentions Donaldson--Sun, singular Yau--Tian--Donaldson, orbifold and synthetic-Ricci theories without precise original-paper references or theorem statements.
\end{itemize}
\end{document}
